\scanpage{15}
% archive foliation: 14
\[
n\longmapsto\operatorname{rg}_K L_n\otimes_k K
\]
est polynomiale pour $n$ grand [le
degré du polynôme \inserted{\unsure{lui-même associé à $L_*$}}
est alors le plus
grand entier $n$ tel que $L^!_n\neq0$, et les
coefficients du polynôme exprimé comme
combinaison linéaire des $P_i(n)=\binom ni$
$(i\geq0)$ sont les $\operatorname{rg}_K L^!_i\otimes_k K$;
les coefficients
dudit polynôme (pour les $P_i$, ou
pour l'écriture habituelle, \unsure{suivant} les
comb. lin. des \unsure{monômes}) sont des
fonctions constructibles en le pt variable
$\mathfrak p\in\operatorname{Spec}k$].

\uline{NB} Si dans e) on suppose $k$ noeth.,
mais non régulier, la condition \inserted{\unsure{énoncée}} de
$\partial$-perfection est tjrs néc, mais pas
suffisante --- il faut une condition
néc et suff. plus subtile, savoir
l'existence d'un $L'_*\longrightarrow L_*$, avec
$L'_*$ \uline{str.} $\partial$-parfait (cf c)) et qui
soit un \deleted{\unsure{isom.}} \inserted{quasi} \unsure{isom.}, i.e. induise un iso
\unsure{sur} les $\pi_i$. (De plus, \unsure{on peut} $L'_*\longrightarrow L_*$
épi.)
\margin{\unsure{Soit $C$ la cat des} \illegible{3}
\unsure{munie des quasi-isomorphismes} \illegible{3}
\unsure{proj. de t.f.} \illegible{2}
$\Lambda^i$ \unsure{donne la} \illegible{2}
\deleted{\unsure{str. $\partial$-parfait}}}

On veut une théorie à \uline{\unsure{isom.}}
\uline{près} des $L_*$ \unsure{(str. $\partial$-ps. coh.)}
\deleted{\unsure{de}}
--- ceci n'est possible bien sûr que
\deleted{\unsure{les ps. cohérents et les $\partial$-parfaits}}
